% Remerciements / Acknowledgements
%
% Grâce aux remerciements, l'auteur attire l'attention du 
% lecteur sur l'aide que certaines personnes lui ont apportée, 
% sur leurs conseils ou sur toute autre forme de contribution 
% lors de la réalisation de son mémoire ou thèse. Le cas 
% échéant, c'est dans cette section que le candidat doit 
% témoigner sa reconnaissance à son directeur de recherche, aux 
% organismes dispensateurs de subventions ou aux entreprises qui
% lui ont accordé des bourses ou des fonds de recherche.

% Through the acknowledgements, the author draws the
% reader's attention to the help that certain people 
% have given them, their advice or any other form of 
% contribution during the completion of the 
% dissertation or thesis. If applicable, it is in 
% this section the candidate should acknowledge the 
% assistance of their advisor, granting agencies or 
% companies that have provided research grants or
% funds.
\ifthenelse{\equal{\Langue}{english}}{
	\chapter*{ACKNOWLEDGEMENTS}\thispagestyle{headings}
	\addcontentsline{toc}{compteur}{ACKNOWLEDGEMENTS}
}{
	\chapter*{REMERCIEMENTS}\thispagestyle{headings}
	\addcontentsline{toc}{compteur}{REMERCIEMENTS}
}
%
% =========================================================================
% NOTE : L'auteur ajoutera ici ses remerciements personnels.
% =========================================================================

\bigskip
\noindent\textbf{Déclaration relative à l'utilisation d'outils d'intelligence artificielle générative}

\medskip
\noindent Conformément au guide \textit{Utilisation des outils d'intelligence artificielle générative} de Polytechnique Montréal (avril 2025), l'auteur déclare avoir recouru aux outils d'intelligence artificielle générative (IAG) suivants dans le cadre de la réalisation de cette thèse~:

\begin{itemize}
    \item \textbf{GitHub Copilot} (Microsoft / OpenAI)~: utilisé pour l'assistance à la rédaction et à la révision du code Python des environnements de simulation (implémentation des protocoles GAPF, CALASH et CASTER-ZT, scripts d'analyse statistique et de visualisation des résultats). Aucune donnée de recherche non publiée, résultat confidentiel ou actif intellectuel valorisable n'a été soumis à cet outil.

    \item \textbf{Claude} (Anthropic)~: utilisé pour la vérification grammaticale et l'amélioration stylistique de passages rédigés en anglais dans les articles soumis à comité de lecture, ainsi que pour la reformulation de sections de la thèse en français. L'ensemble des idées, contributions scientifiques originales et résultats numériques demeurent entièrement l'œuvre de l'auteur.
\end{itemize}

\noindent L'utilisation de ces outils ne porte sur aucun actif intellectuel valorisable et se limite aux tâches d'assistance à la programmation et à l'édition linguistique, conformément aux exceptions prévues par la politique en vigueur. Aucune information personnelle, confidentielle ou stratégique de Polytechnique Montréal n'a été communiquée à ces outils.
